% Vietnamese translation for OI Public Library
% Source-File: 比赛 CONTEST/2021“MINIEYE杯”中国大学生算法设计超级联赛/2021“MINIEYE杯”中国大学生算法设计超级联赛（1）/Necklace of Beads/solution.docx
\documentclass[11pt]{article}
\usepackage{oipl-vn}

\newcommand{\ans}{\operatorname{ans}}

\title{Necklace of Beads: lời giải}
\author{Wo Wei Shu Xue Kuang}
\date{20 tháng 7, 2021}

\begin{document}
\maketitle

\origtitle{Necklace of Beads solution}
\sourcepath{Contest / 2021 MINIEYE Cup, Multi-University Training Contest 1 / Necklace of Beads / solution.docx}

\section{Ý tưởng}

Vì phải xét các phương án trùng nhau sau khi quay vòng, ta dùng bổ đề
Burnside.

Ký hiệu \(g(n,m)\) là số cách tô một vòng kích thước \(n\), trong đó số hạt
màu xanh lá là \(m\), và mọi điều kiện của bài toán đều được thỏa mãn. Khi
quay vòng đi \(i\) vị trí, số chu trình của hoán vị là \(\gcd(i,n)\). Theo
Burnside, đáp án có dạng
\[
  \ans
  =
  \sum_{i=0}^{n-1}
  \sum_{j=0}^{\left\lfloor k\gcd(i,n)/n\right\rfloor}
  g(\gcd(i,n),j).
\]

Lý do có cận trên của \(j\) là một vị trí trong vòng sau khi rút gọn theo
chu trình tương ứng với \(n/\gcd(i,n)\) vị trí trên vòng ban đầu.

\section{\texorpdfstring{Tính \(g(n,m)\)}{Tính g(n,m)}}

Vì số lượng của một màu bị hạn chế, ta có thể trước hết chọn vị trí của màu
đó. Giả sử đã đặt hạt xanh lá lên \(m\) vị trí đôi một không kề nhau trên
vòng. Trong các khoảng trống giữa chúng, các màu còn lại có đúng hai lựa
chọn độc lập, nên
\[
  g(n,m)=2^m h(n,m),
\]
trong đó \(h(n,m)\) là số cách chọn \(m\) vị trí không kề nhau trên một vòng
kích thước \(n\).

Xét hai trường hợp:
\begin{itemize}
  \item Nếu không chọn vị trí \(n\), số cách là \(\binom{n-m}{m}\).
  \item Nếu chọn vị trí \(n\), hai vị trí \(1\) và \(n-1\) đều không được
  chọn, số cách là \(\binom{n-m-1}{m-1}\).
\end{itemize}
Vì vậy
\[
  g(n,m)=2^m\left(\binom{n-m}{m}+\binom{n-m-1}{m-1}\right),
\]
với trường hợp \(g(n,0)\) cần được xét riêng.

\section{Rút gọn tổng theo ước}

Thay công thức Burnside ở trên, rồi nhóm các giá trị có cùng
\(\gcd(i,n)=d\):
\[
\begin{aligned}
  \ans
  &=
  \sum_{i=0}^{n-1}
  \sum_{j=0}^{\left\lfloor k\gcd(i,n)/n\right\rfloor}
  g(\gcd(i,n),j) \\
  &=
  \sum_{d=1}^{n}
  \left(\sum_i [\gcd(n,i)=d]\right)
  \sum_{j=0}^{\left\lfloor kd/n\right\rfloor}g(d,j) \\
  &=
  \sum_{d\mid n}
  \varphi\!\left(\frac{n}{d}\right)
  \sum_{j=0}^{\left\lfloor kd/n\right\rfloor}g(d,j).
\end{aligned}
\]
Nếu đặt
\[
  F(d)=\sum_{j=0}^{\left\lfloor kd/n\right\rfloor}g(d,j),
\]
thì
\[
  \ans=\sum_{d\mid n}\varphi\!\left(\frac{n}{d}\right)F(d).
\]

\section{Cài đặt}

Tiền xử lý giai thừa và nghịch đảo giai thừa để tính nhanh tổ hợp. Sau đó
liệt kê các ước \(d\mid n\), với mỗi \(d\) tính \(F(d)\) bằng một vòng lặp
theo \(j\). Ghi chú trong nguồn nói rằng với mỗi \(d\), tính trực tiếp trong
\(O(n)\) là đủ.

\section{Ghi chú từ nguồn}

Do máy trên Aliyun chạy rất nhanh, một số đội có tối ưu chưa thật hoàn chỉnh
vẫn vượt qua được bài này.

\end{document}
